This thesis set out to connect a detailed picture of the U.S. used fuel inventory
to a model of an accelerator driven burner, so that the transmutation of minor
actinides could be evaluated as a national fuel cycle rather than as an isolated
reactor calculation. Two cyclus archetypes were developed to that end.
\texttt{us\_inventory} represents the inventory at assembly resolution, supplying
material by a choice of selection policies while preserving each assembly's
composition and provenance; \texttt{ads} models an accelerator driven
minor actinide burner from precomputed Serpent depletion cases, transforming a
loaded core into its irradiated composition and closing the actinide mass balance
each cycle. Together they form a complete transmutation fuel cycle when combined
with standard separations, mixing, and disposal facilities.

Exercising this fuel cycle across eight scenarios produced two principal
findings. First, minor actinide transmutation is limited by the availability
of its feed rather than by ADS capacity. The accessible minor actinide
inventory is approximately \(103.8~\mathrm{t}\), fewer than four fresh
loadings of the \(27~\mathrm{t}\) ADS~1 core. ADS~1 therefore stalls after
three cycles under once through operation and five cycles with recycling.
ADS~2 avoids this limitation only by coburning plutonium, thereby changing
the mission from dedicated minor actinide destruction to broader transuranic
destruction.

Second, the direct input--output comparisons show that neither ADS
configuration eliminates its loaded actinides in one irradiation pass.
ADS~1 reduces the actinide mass from \(27.00~\mathrm{t}\) to approximately
\(15.84~\mathrm{t}\), corresponding to \(11.16~\mathrm{t}\), or
\(41.3\%\), of net actinide destruction per cycle. ADS~2 reduces the
actinide mass from \(4.418~\mathrm{t}\) to approximately
\(2.231~\mathrm{t}\), corresponding to \(2.187~\mathrm{t}\), or
\(49.5\%\), of net destruction per cycle. The remaining actinides must be
recovered and irradiated again if greater destruction is required. At the
full fuel cycle scale, the 100 year ADS~1 recycling scenario destroys
\(55.8~\mathrm{t}\) of the initial \(103.8~\mathrm{t}\) accessible
minor actinide inventory, while the once through case destroys
\(33.5~\mathrm{t}\).

These results show that accelerator driven transmutation is not a stand alone
solution to the spent fuel problem. It is a downstream component of a much
larger reprocessing and recycle system. Before the burner can provide any
benefit, essentially the entire selected spent fuel inventory must be
retrieved and processed to recover the comparatively small plutonium and
minor actinide fractions. The simulations assume that these materials can be
recovered into separate streams at high efficiency; they do not demonstrate
the technical, economic, safeguards, licensing, or policy feasibility of that
infrastructure. The difficulty of establishing and operating the required
reprocessing system may therefore be at least as important as the difficulty
of designing the ADS itself.

Partitioning also changes the composition and management of the waste streams
before any ADS irradiation occurs. Separating uranium, fission products,
plutonium, and minor actinides may itself alter decay heat, radiotoxicity, waste
form requirements, and repository performance, depending on the disposition
of each stream. Because this thesis does not include a reprocessing only
reference scenario or a downstream radiotoxicity calculation, it cannot
determine how much of the potential long term waste benefit would result from
partitioning alone and how much would result from the additional actinide
destruction in the ADS. The contribution quantified here is therefore the
incremental net destruction performed by the burner after separation:
\(41.3\%\) per pass for ADS~1 and \(49.5\%\) for ADS~2.

Several extensions follow directly from this work. The most immediate is to
couple the mass-conserving waste histories to a decay and dose calculation so
that the reductions in long term radiotoxicity and decay heat, which motivate
minor actinide removal, can be quantified rather than inferred. That analysis
should include a reprocessing only reference scenario in which uranium,
plutonium, minor actinides, and fission products are partitioned but not
irradiated in an ADS. Comparing that case with both the unreprocessed inventory
and the ADS scenarios would distinguish the repository benefit of separation
itself from the incremental benefit of transmutation.

A second extension is to use the assembly-level provenance recorded by
\texttt{us\_inventory}, which associates each supplied assembly with its source
facility and coordinates, as the input to a transportation and siting analysis
of where retrieval and processing would occur. A third is to develop an
economic model of the reprocessing and handling burden implied by the results,
which this study describes but does not price. Finally, the reduced order
\texttt{ads} model could be broadened with a larger library of depletion cases,
and the effects of inventory selection policies could be examined within the
fully coupled simulation rather than only against the static inventory. In
each case, the archetypes and results developed here provide a foundation for
these subsequent analyses.